% preview-header.tex — LaTeX preamble that encodes the journal's format rules.
%
% Used by paper/preview.ps1 (pandoc -> XeLaTeX -> PDF) to render the drafts the
% way the printed article will look, WITHOUT going through Word.  It is the
% LaTeX counterpart of journal/templates/reference.docx: reference.docx carries
% the format for the .docx that gets submitted, this file carries it for the
% on-screen preview.  Every rule below is traceable to journal/FORMAT_GUIDE.md.
%
% Page geometry, font family and base size are passed on the pandoc command
% line in preview.ps1, not set here.

% --- s.1  จำนวนคอลัมน์: 2 คอลัมน์ ระยะห่าง 0.5 ซม. -------------------------
% (column WIDTH follows from the page geometry; see the note in preview.ps1)
\setlength{\columnsep}{0.5cm}

% --- s.1  ระยะระหว่างบรรทัด: บรรทัดเดี่ยว (Single Space) --------------------
% TH Sarabun New declares hhea asc/desc/gap = 844/-457/30 per 1000 upem
% = 1.331 em, so Word's "single" spacing leads at 14 x 1.331 = 18.6 pt.  The
% 12 pt class leads at 14.5 pt, hence 18.6/14.5 = 1.285.  Without this the
% preview is ~25% denser than the .docx and under-reports the page count.
\linespread{1.285}

% --- Thai line breaking ----------------------------------------------------
% XeTeX does not wrap Thai at all without a Thai locale; lines would simply
% overflow the column.  The break glue must stay at 0pt with NO stretch: any
% "plus" component is stretchable glue sitting at every Thai word boundary, and
% justification pulls it open into visible gaps in the middle of Thai text,
% which has no inter-word spaces to begin with.  Word does this correctly with
% Thai Distribute (set in reference.docx); XeTeX cannot, so the preview is left
% slightly ragged rather than wrongly spaced.
\XeTeXlinebreaklocale "th"
\XeTeXlinebreakskip = 0pt plus 0.3pt
% Let TeX try much harder before it resorts to a badly stretched line.  Without
% this it hits a bad break, then has to open the few available gaps very wide -
% which is what produced the visible holes inside Thai text.
\tolerance = 4000
\emergencystretch = 3em

\usepackage{amsmath}

% --- ตาราง ------------------------------------------------------------------
% pandoc แปลงทุกตารางเป็น longtable ซึ่งใช้ในโหมดสองคอลัมน์ไม่ได้ (LaTeX แจ้ง
% "longtable not in 1-column mode") ตารางกว้างในบทความสองคอลัมน์ควรวางเป็น
% float เต็มความกว้างหน้าอยู่แล้ว จึงเปลี่ยนเส้นทาง longtable ไปเป็น table*
% พร้อม tabular ธรรมดา และปิดคำสั่งหัว/ท้ายที่มีเฉพาะใน longtable
% pandoc วาง \bottomrule ไว้ในตำแหน่งของ "foot" คือก่อนแถวข้อมูล ซึ่งถูกต้อง
% สำหรับ longtable แต่ผิดตำแหน่งใน tabular จึงปิดมันไว้ระหว่างตารางแล้วค่อย
% ออกเส้นล่างเองตอนจบ environment
\usepackage{array,booktabs}
\let\origbottomrule\bottomrule
\newenvironment{pandoclongtable}[2][]
  {\begin{table*}[tbp]\centering\def\bottomrule\noalign##1{}\begin{tabular}{#2}}
  {\origbottomrule\end{tabular}\end{table*}}
\let\longtable\pandoclongtable
\let\endlongtable\endpandoclongtable
\let\endhead\relax
\let\endfirsthead\relax
\let\endfoot\relax
\let\endlastfoot\relax

% --- ย่อหน้า ---------------------------------------------------------------
% The drafts open every paragraph with &emsp;.  preview.ps1 strips that entity
% (XeLaTeX has no glyph for U+2003 in TH Sarabun New) and the indent comes from
% here instead; indentfirst extends it to the first paragraph after a heading,
% which is what the drafts do.
\setlength{\parindent}{1.5em}
\usepackage{indentfirst}
\setlength{\parskip}{0pt}

% --- s.3  ขนาดและรูปแบบหัวข้อ ----------------------------------------------
% All heading levels sit at body size (14 pt); main and sub headings are bold,
% sub-sub headings are regular.
% --- s.5  เว้น 1 บรรทัดก่อนหัวข้อหลักและหัวข้อรอง (หัวข้อย่อยไม่ต้องเว้น) ---
% pandoc maps markdown ## -> \subsection, ### -> \subsubsection,
% #### -> \paragraph for the article class, so the levels are shifted by one
% relative to the names in FORMAT_GUIDE.md.
\usepackage{titlesec}
\titleformat{\section}{\bfseries}{}{0pt}{}
\titleformat{\subsection}{\bfseries}{}{0pt}{}
\titleformat{\subsubsection}{\normalfont}{}{0pt}{}
\titleformat{\paragraph}[block]{\normalfont}{}{0pt}{}
\titlespacing*{\section}{0pt}{1\baselineskip}{0pt}
\titlespacing*{\subsection}{0pt}{1\baselineskip}{0pt}
\titlespacing*{\subsubsection}{0pt}{0pt}{0pt}
\titlespacing*{\paragraph}{0pt}{0pt}{0pt}

% หัวข้อถูกใส่เลขไว้ในต้นฉบับเองแล้ว (2.1, 2.2.3, ...) — อย่าให้ LaTeX ใส่ซ้ำ
\setcounter{secnumdepth}{0}
